%%%%%%%%%%%%%%%%%%%%%%%%%%%%%%%%%%%%%%%%%%%%%%%%%%%%%%%%%%%%%%%%%%%%%%%%%%%%%%%%
\chapter{User Manual}
\label{chap:user_manual}

This chapter provides a complete user manual for the RANPilotAI platform. It covers system requirements, installation, launching the Telecom Hub, and step-by-step operating instructions for each of the four analytical modules.

%%%%%%%%%%%%%%%%%%%%%%%%%%%%%%%%%%%%%%%%%%%%%%%%%%%%%%%%%%%%%%%%%%%%%%%%%%%%%%%%
\section{System Requirements}

\begin{itemize}
    \item \textbf{Operating System:} Windows 10/11, Linux, or macOS
    \item \textbf{Python:} Version 3.11 or higher
    \item \textbf{RAM:} 8\,GB minimum (16\,GB recommended for large KPI datasets)
    \item \textbf{Disk Space:} At least 5\,GB free (additional space required for 3GPP PDF files and the local vector database)
    \item \textbf{Internet Connection:} Required for the RAG Assistant (LLM API calls via OpenRouter); all other modules operate fully offline
    \item \textbf{Web Browser:} Any modern browser (Chrome, Firefox, Edge) for the Streamlit interface
\end{itemize}

%%%%%%%%%%%%%%%%%%%%%%%%%%%%%%%%%%%%%%%%%%%%%%%%%%%%%%%%%%%%%%%%%%%%%%%%%%%%%%%%
\section{Installation}

\subsection{Step 1 — Install Python}

Ensure Python 3.11 or higher is installed. Verify by running:

\begin{verbatim}
python --version
\end{verbatim}

\subsection{Step 2 — Install Dependencies}

From the project root directory, install all required packages with a single command:

\begin{verbatim}
pip install -r requirements.txt
\end{verbatim}

This installs all dependencies for all four tools, including Streamlit, Pandas, NumPy, SciPy, XGBoost, Statsmodels, Sentence Transformers, Qdrant, PyMuPDF, and python-docx.

\subsection{Step 3 — Configure the RAG Assistant (First Time Only)}

If you intend to use the Telecom RAG Assistant, an OpenRouter API key is required:

\begin{enumerate}
    \item Copy the environment template:
    \begin{verbatim}
cp RAG_tool/.env.example RAG_tool/.env
    \end{verbatim}
    \item Open \texttt{RAG\_tool/.env} and enter your OpenRouter API key:
    \begin{verbatim}
OPENROUTER_API_KEY=your_api_key_here
    \end{verbatim}
    \item Place your 3GPP specification PDF files in \texttt{RAG\_tool/data/} using the naming convention:
    \begin{verbatim}
TS_<NUMBER>_Rel_<RELEASE>.pdf
    \end{verbatim}
    For example: \texttt{TS\_36.331\_Rel\_17.pdf}
    \item Run the ingestion pipeline to index the PDFs into the local vector database:
    \begin{verbatim}
python -m app.ingest --all
    \end{verbatim}
    Note: The first run downloads the BGE-large embedding model ($\approx$1.3\,GB).
\end{enumerate}

\subsection{Step 4 — Configure the Optimization Action Analyzer (If Required)}

The Optimization Action Analyzer requires a running Dolt SQL server with two tables: \texttt{network\_kpis} and \texttt{network\_configs}. Install Dolt from \url{https://docs.dolthub.com/introduction/installation} and start the server on the default port before launching the platform.

%%%%%%%%%%%%%%%%%%%%%%%%%%%%%%%%%%%%%%%%%%%%%%%%%%%%%%%%%%%%%%%%%%%%%%%%%%%%%%%%
\section{Launching the Platform}

The entire platform is launched through a single entry point from the project root directory:

\begin{verbatim}
streamlit run telecom_hub.py
\end{verbatim}

After a few seconds, Streamlit will display a local URL (typically \texttt{http://localhost:8501}). Open this URL in your web browser. The Telecom Hub landing page will appear, displaying cards for all four analytical tools.

%%%%%%%%%%%%%%%%%%%%%%%%%%%%%%%%%%%%%%%%%%%%%%%%%%%%%%%%%%%%%%%%%%%%%%%%%%%%%%%%
\section{Navigating the Telecom Hub}

The Telecom Hub interface consists of two areas:

\begin{itemize}
    \item \textbf{Left Sidebar:} A dropdown menu listing all four tools. Select any tool to load it in the main area.
    \item \textbf{Main Content Area:} Displays the landing page or the selected tool's full interface.
\end{itemize}

To switch between tools, use the sidebar dropdown at any time. No page reload or server restart is required.

%%%%%%%%%%%%%%%%%%%%%%%%%%%%%%%%%%%%%%%%%%%%%%%%%%%%%%%%%%%%%%%%%%%%%%%%%%%%%%%%
\section{Tool 1 — KPI Degradation Analyzer}

\subsection{Purpose}

Automatically detects performance degradation in LTE and 5G NR cells by comparing a recent KPI window against a historical baseline, ranking root causes, and generating exportable reports.

\subsection{Input Data Format}

The tool accepts Excel files (\texttt{.xlsx} or \texttt{.xls}) containing daily KPI exports. The file must include:

\begin{itemize}
    \item An \texttt{eNodeB Name} column and a \texttt{Cell Name} (or \texttt{LocalCell Id}) column for cell identification
    \item A \texttt{Date} column in a standard date format
    \item One or more KPI counter columns (column names are matched flexibly using fuzzy matching)
\end{itemize}

\subsection{Step-by-Step Usage}

\begin{enumerate}
    \item Select \textbf{KPI Degradation Analyzer} from the sidebar dropdown.
    \item Click \textbf{Upload File} and select your Excel KPI export.
    \item Select the appropriate worksheet from the sheet dropdown.
    \item Choose the \textbf{KPI} to analyze from the KPI selector.
    \item Set the \textbf{Comparison Days} (the number of recent days to compare against the baseline; default is 4).
    \item Set the \textbf{Degradation Threshold} (the minimum degradation percentage to flag a cell; default is 30\%).
    \item Select the \textbf{Baseline Mode}:
        \begin{itemize}
            \item \textit{Last Week} — compares the recent window against the same days in the previous week. Best for detecting sudden incidents.
            \item \textit{4-Week Rolling Median} — uses the median of the same weekdays over the past four weeks. Best for smoothing recurring weekly patterns.
        \end{itemize}
    \item Click \textbf{Run Selected KPI} to analyze one KPI, or \textbf{Analyze All KPIs} to process all 13 categories at once.
    \item Review the results table displaying degraded cells, degradation percentages, root cause categories, and recommended actions.
    \item Use the \textbf{Charts} tab to view bar charts of degraded cells per KPI and root cause distribution.
    \item Use the \textbf{Trends} tab to view before/after line chart overlays.
    \item Export results using the \textbf{Save CSV} button or generate a formatted Word report with \textbf{Generate Report}.
\end{enumerate}

\subsection{Understanding the Output}

\begin{itemize}
    \item \textbf{rf\_severity}: Operational RF impact label — Normal, Medium, High, or Critical.
    \item \textbf{rca\_pattern}: Root cause pattern — Outage, Congestion, Radio Quality, Coverage, Interference, Mobility, Demand, or Unknown.
    \item \textbf{main\_root\_cause\_category}: The highest-scored root cause category.
    \item \textbf{main\_recommended\_action}: The engineering action recommended for the top cause.
    \item \textbf{analysis\_confidence}: Confidence label based on statistical significance, data completeness, and severity.
\end{itemize}

%%%%%%%%%%%%%%%%%%%%%%%%%%%%%%%%%%%%%%%%%%%%%%%%%%%%%%%%%%%%%%%%%%%%%%%%%%%%%%%%
\section{Tool 2 — KPI Forecasting Engine}

\subsection{Purpose}

Forecasts future KPI trends using machine learning and statistical time-series models, enabling engineers to anticipate degradation before it occurs operationally.

\subsection{Input Data Format}

A CSV file with:
\begin{itemize}
    \item A \texttt{Date} column in \texttt{YYYY-MM-DD} format
    \item A \texttt{Cell Name} column
    \item One column per KPI using the raw counter names defined in \texttt{config.py}
\end{itemize}

\subsection{Step-by-Step Usage}

\begin{enumerate}
    \item Select \textbf{KPI Forecaster} from the sidebar dropdown.
    \item Upload your cleaned KPI CSV file when prompted.
    \item From the sidebar, choose a \textbf{Forecasting Method}: XGBoost, Holt-Winters, or Compare Both.
    \item Select the \textbf{KPI to Forecast} from the dropdown (12 KPIs supported).
    \item Select a \textbf{Cell} from the dropdown (sorted by weekly seasonality strength, strongest first).
    \item Set the \textbf{Hold-out Test Days} (the number of most recent days used for backtesting; default is 4).
    \item Optionally enable:
        \begin{itemize}
            \item \textit{Show Next-7-Day Future Forecast} — extends the prediction beyond the test window.
            \item \textit{Flag Seasonality Strength} — displays STL-based seasonality badges.
            \item \textit{Compare vs Naive Baseline} — benchmarks models against naive and weekly-mean baselines.
        \end{itemize}
    \item Adjust XGBoost or Holt-Winters hyperparameters from the sidebar if required.
    \item Review the forecast chart, accuracy metrics (MAE, RMSE, MAPE), predictions table, and the 7-day forward forecast table.
    \item Expand the \textbf{Feature Importance} panel to inspect the top-15 XGBoost predictors.
    \item Expand the \textbf{Residual Analysis} panel for diagnostic statistics (Durbin-Watson, Ljung-Box).
    \item Use the \textbf{LLM-Ready Report} panel to generate a structured summary suitable for pasting into an LLM for further optimization reasoning.
\end{enumerate}

%%%%%%%%%%%%%%%%%%%%%%%%%%%%%%%%%%%%%%%%%%%%%%%%%%%%%%%%%%%%%%%%%%%%%%%%%%%%%%%%
\section{Tool 3 — Optimization Action Analyzer}

\subsection{Purpose}

Correlates RF configuration changes (antenna tilt adjustments, power parameter updates) recorded in a Dolt version-controlled database with before/after KPI impact, enabling engineers to quantitatively evaluate the effect of optimization actions.

\subsection{Prerequisites}

\begin{itemize}
    \item Dolt installed and a Dolt SQL server running on the default port (3306)
    \item A Dolt database containing \texttt{network\_kpis} and \texttt{network\_configs} tables
    \item Configuration changes committed to the Dolt database (each commit corresponds to one optimization action)
\end{itemize}

\subsection{Step-by-Step Usage}

\begin{enumerate}
    \item Select \textbf{Optimization Action Analyzer} from the sidebar dropdown.
    \item The tool connects to the Dolt server automatically and reads the commit log.
    \item Select an \textbf{Action Hunk} (a date grouping of one or more Dolt commits) from the timeline to analyze.
    \item For each KPI detected in the database, set the \textbf{polarity}:
        \begin{itemize}
            \item \textit{Higher is Better} (e.g., Throughput, Setup Success Rate)
            \item \textit{Lower is Better} (e.g., Drop Rate, Block Rate)
        \end{itemize}
    \item Review the results:
        \begin{itemize}
            \item \textbf{Action Timeline} — a visual chart of all Dolt commits.
            \item \textbf{Before/After Line Chart} — overlays the same weekday before and after the action.
            \item \textbf{Delta Bar Chart} — shows the absolute KPI change per time interval, color-coded green (improved) or red (degraded).
            \item \textbf{KPI Summary Table} — structured before/after comparison with absolute delta, percentage change, and status.
        \end{itemize}
\end{enumerate}

%%%%%%%%%%%%%%%%%%%%%%%%%%%%%%%%%%%%%%%%%%%%%%%%%%%%%%%%%%%%%%%%%%%%%%%%%%%%%%%%
\section{Tool 4 — Telecom RAG Assistant}

\subsection{Purpose}

Answers engineering questions about LTE and 5G NR standards in natural language, with every response grounded in and cited from locally indexed 3GPP Technical Specification PDFs.

\subsection{Prerequisites}

Ensure that the RAG Assistant was configured as described in Step~3 of the Installation section: the \texttt{.env} file must contain a valid OpenRouter API key, 3GPP PDFs must be placed in \texttt{RAG\_tool/data/}, and the ingestion pipeline must have been run at least once.

\subsection{Step-by-Step Usage}

\begin{enumerate}
    \item Select \textbf{Telecom RAG Assistant} from the sidebar dropdown.
    \item Optionally use the \textbf{Specification Filter} in the sidebar to restrict answers to a specific TS number (e.g., \texttt{36.331}) or 3GPP release (e.g., \texttt{17}).
    \item Type your engineering question in natural language in the chat input box. Examples:
        \begin{itemize}
            \item \textit{``What are the trigger conditions for Event A3 in LTE RRC?''}
            \item \textit{``What is the purpose of the RACH procedure in 5G NR?''}
            \item \textit{``Explain the difference between RLC AM and UM modes.''}
        \end{itemize}
    \item Press Enter or click Send.
    \item The assistant will display a structured technical response followed by expandable citation panels showing the exact TS number, release version, section number, and page number for each retrieved chunk.
    \item If the answer cannot be found in the indexed specifications, the assistant will explicitly state so rather than generating an unverified response.
\end{enumerate}

\subsection{Performance Tuning}

The retrieval behaviour can be adjusted via the \texttt{RAG\_tool/.env} file:

\begin{itemize}
    \item \texttt{TOP\_K} (default: 5) — number of specification chunks supplied to the language model.
    \item \texttt{FETCH\_MULTIPLIER} (default: 4) — multiplier for chunks retrieved from Qdrant before re-ranking.
    \item \texttt{CHUNK\_SIZE} (default: 800 tokens) — maximum size of each indexed chunk.
    \item \texttt{DEBUG\_MODE} (default: false) — enables pipeline stage timing and prompt statistics in the sidebar.
\end{itemize}

%%%%%%%%%%%%%%%%%%%%%%%%%%%%%%%%%%%%%%%%%%%%%%%%%%%%%%%%%%%%%%%%%%%%%%%%%%%%%%%%
\section{Troubleshooting}

\begin{description}
    \item[The platform does not start.] Verify that Python 3.11 or higher is installed and that all dependencies were installed with \texttt{pip install -r requirements.txt}. Check that \texttt{telecom\_hub.py} is executed from the project root directory.

    \item[The KPI Degradation Analyzer cannot find KPI columns.] Ensure your Excel file includes column headers that partially match the expected counter names. The tool uses fuzzy matching, but column names must be recognizable. Refer to the KPI configuration reference in \texttt{KPI\_and\_its\_related\_counters.md}.

    \item[The Forecaster produces a flat or unstable forecast.] This may occur when the data window is shorter than 21 days, causing lag features to be sparsely populated. Reduce \texttt{max\_depth} and \texttt{n\_estimators} in the XGBoost sidebar, or switch to the Holt-Winters model.

    \item[The Optimization Action Analyzer cannot connect to Dolt.] Verify that the Dolt SQL server is running (\texttt{dolt sql-server}) and that the \texttt{network\_kpis} and \texttt{network\_configs} tables exist in the active database.

    \item[The RAG Assistant returns ``I cannot find this in the supplied specifications.''] The queried concept may not be present in the indexed PDFs. Verify that the relevant 3GPP specification has been placed in \texttt{RAG\_tool/data/} and re-run the ingestion pipeline (\texttt{python -m app.ingest --all}).

    \item[The RAG Assistant is slow to respond.] Response latency depends on the OpenRouter API. Enable \texttt{DEBUG\_MODE=true} in \texttt{.env} to display per-stage pipeline timing in the sidebar and identify bottlenecks.
\end{description}
